四矢量表示给出时空坐标在 Lorentz 变换下的线性混合规律；自旋 $1/2$ 粒子的量子态还必须容纳一种四矢量表示看不见的性质：空间转动 $2\pi$ 后，所有四矢量已经回到原值，但量子态仍可以多出一个负号。因此仅用 $4\times4$ Lorentz 矩阵作用在四矢量上，还不足以描述旋量怎样随转动和推动变化。因此需要构造一组 $2\times2$ 复矩阵，使每个保向、正时向 Lorentz 变换对应其中两个相差整体负号的矩阵。这种“两个矩阵对应同一个 Lorentz 变换”的二对一关系称为\emph{双覆盖}，相应矩阵群就是 $SL(2,\mathbb C)$。


把实四矢量 $x^\mu=(x^0,\mathbf x)$ 与厄米 $2\times2$ 矩阵一一对应：
\begin{equation}
 X(x)\equiv x^0\id(2)+\mathbf x\cdot\boldsymbol\sigma_{\rm P}
 =\begin{pmatrix}
 x^0+x^3 & x^1-\ii x^2\\
 x^1+\ii x^2 & x^0-x^3
 \end{pmatrix}.
 \label{eq:sl2c-hermitian-map-dp8}
\end{equation}
这个对应之所以有用，是因为矩阵行列式恰好等于 Minkowski 范数：
\begin{equation}
 \boxed{\det X(x)=x_\mu x^\mu.}
 \label{eq:sl2c-det-minkowski-dp8}
\end{equation}
因此，只要找到保持 $\det X$ 的矩阵变换，就自动得到保持时空间隔的变换。

对任意 $A\in SL(2,\mathbb C)$，定义
\begin{equation}
 X' = A X A^\dagger.
 \label{eq:sl2c-action-X-dp8}
\end{equation}
$X'$ 仍为厄米矩阵，因而唯一对应某个实四矢量 $x'^\mu$；同时 $\det A=1$ 保证 $\det X'=\det X$。所以每个 $A$ 都诱导一个保向、正时向 Lorentz 变换。另一方面，$A$ 与 $-A$ 对 $X$ 的作用完全相同，于是得到二对一群同态
\begin{equation}
 \boxed{SL(2,\mathbb C)\longrightarrow SO^+(1,3),\qquad A\sim -A.}
 \label{eq:sl2c-double-cover-dp8}
\end{equation}
这就是所谓“双覆盖”：同一个 Lorentz 变换对应 $SL(2,\mathbb C)$ 中相差负号的两个元素。

这里必须区分实 Lie 代数与其复化。实 Lorentz 代数 $\mathfrak{so}(1,3)$ 本身是实简单 Lie 代数，并不等于两份 $\mathfrak{su}(2)$ 的实直和。为了分类有限维复表示，把生成元允许作复线性组合；此时
\begin{equation*}
 \mathfrak{so}(1,3)_\mathbb C
 \simeq \mathfrak{sl}(2,\mathbb C)\oplus\mathfrak{sl}(2,\mathbb C),
\end{equation*}
等价地可用 $\mathsf J_i^{(\pm)}=(J_i\pm\ii K_i)/2$ 得到两套彼此对易、具有 $\mathfrak{su}(2)$ 型结构常数的复子代数。有限维复不可约表示因此由两个角动量型标签 $(j_L,j_R)$ 标记，其复维数为
\begin{equation}
 \boxed{\dim(j_L,j_R)=(2j_L+1)(2j_R+1).}
 \label{eq:lorentz-irrep-dimension-dp8}
\end{equation}
最常用的表示由此统一写成
\begin{equation}
 \begin{array}{ccl}
 (0,0) &:& \text{标量},\\[2pt]
 (\tfrac12,0) &:& \text{左手 Weyl 旋量},\\[2pt]
 (0,\tfrac12) &:& \text{右手 Weyl 旋量},\\[2pt]
 (\tfrac12,\tfrac12) &:& \text{四矢量},\\[2pt]
 (1,0)\oplus(0,1) &:& \text{反对称二阶张量},\\[2pt]
 (\tfrac12,0)\oplus(0,\tfrac12) &:& \text{Dirac 旋量}.
 \end{array}
 \label{eq:lorentz-irrep-dictionary-dp8}
\end{equation}
这里的 $(j_L,j_R)$ 描述的是\emph{局域场怎样变换}；它与后面由 Poincar\'e 不可约表示得到的单粒子自旋标签相关，但不是同一个概念。


表示的张量积描述复合 Lorentz 指标怎样重新组合，而不是把两个粒子的自旋数简单相加。左右 Weyl 指标各取一次时，维数 $2\times2=4$，并且
\begin{equation}
 \boxed{
 \left(\frac12,0\right)\otimes\left(0,\frac12\right)
 =\left(\frac12,\frac12\right),
 \qquad
 V_{\alpha\dot\beta}=V_\mu(\sigma^\mu)_{\alpha\dot\beta}.}
 \label{eq:weyl-vector-product-r68}
\end{equation}
所以一个带左右各一个 Weyl 指标的 $2\times2$ 厄米对象正好按四矢量表示变换。若两个指标具有相同手征性，则
\begin{equation}
 \boxed{
 \left(\frac12,0\right)\otimes\left(\frac12,0\right)
 =(0,0)\oplus(1,0),
 \qquad
 \left(0,\frac12\right)\otimes\left(0,\frac12\right)
 =(0,0)\oplus(0,1).}
 \label{eq:weyl-same-chirality-product-r68}
\end{equation}
四维反对称二阶张量的六个分量因此分成两个三维手征部分，写成
\begin{equation}
 \boxed{F^{\mu\nu}\in(1,0)\oplus(0,1).}
 \label{eq:twoform-chiral-decomposition-r68}
\end{equation}
这就是电磁场强可以进一步组织成自对偶与反自对偶组合的表示论来源。

\begin{figure}[htbp]
 \centering
 \resizebox{0.93\textwidth}{!}{\begin{tikzpicture}[x=1cm,y=1cm,>=Stealth,every node/.style={font=\small}]
% Block-logic grid: one row per layer, orthogonal arrows only.
%   row 1 (top) : SO -- SL            (covering map)
%   row 2       : L , R               (chiral reps of SL(2,C))
%   row 3       : build               (combine L and R)
%   row 4 (bot) : D , V , F           (composite reps)
% Same-layer boxes share one minimum width; arrows are horizontal/vertical or
% L-shaped via -| .  All labels use blocklabel (= line-label, white backing).

% ---- row 1: covering map ---------------------------------------------
\node[blockbox,minimum width=48mm] (SO) at (-7.0,3.0)
  {$SO^+(1,3)$\\proper orthochronous\\Lorentz group};
\node[blockbox,minimum width=48mm] (SL) at (0,3.0)
  {$SL(2,\mathbb C)$};
\draw[layerarrow] (SL.west) -- node[blocklabel,above=4pt]{2:1 covering map} (SO.east);

% ---- row 2: chiral representations -----------------------------------
\node[blockbox,minimum width=46mm] (L) at (-3.6,0.0)
  {$(\tfrac12,0)$\\left-handed Weyl spinor $L$};
\node[blockbox,minimum width=46mm] (R) at (3.6,0.0)
  {$(0,\tfrac12)$\\right-handed Weyl spinor $R$};
% L-shaped: leave SL straight down, run out, then drop into each chiral box.
\draw[layerarrow] (SL.south) -| node[blocklabel,pos=0.80,left]{representation} (L.north);
\draw[layerarrow] (SL.south) -| node[blocklabel,pos=0.80,right]{representation} (R.north);

% ---- row 3: build node ------------------------------------------------
% Separate operation node prevents arrows from falsely suggesting conversion.
\node[blockbox,minimum width=58mm] (build) at (0,-2.2)
  {build new Lorentz representations\\from $L$ and $R$};
\draw[layerarrow] (L.south) -| (build.north west);
\draw[layerarrow] (R.south) -| (build.north east);

% ---- row 4: composite reps --------------------------------------------
\node[blockbox,minimum width=50mm] (D) at (-5.0,-4.6)
  {$L\oplus R$\\Dirac spinor};
\node[blockbox,minimum width=50mm] (V) at (0,-4.6)
  {$L\otimes R$\\four-vector $A^\mu$};
\node[blockbox,minimum width=50mm] (F) at (5.0,-4.6)
  {$\mathrm{Sym}^2L\oplus\mathrm{Sym}^2R$\\antisymmetric tensor $F^{\mu\nu}$};
\draw[layerarrow] (build.south west) -| node[blocklabel,pos=0.80,left]{$\oplus$} (D.north);
\draw[layerarrow] (build.south) -- node[blocklabel,right]{$\otimes$} (V.north);
\draw[layerarrow] (build.south east) -| node[blocklabel,pos=0.80,right]{$\mathrm{Sym}^2$} (F.north);
\end{tikzpicture}
}
 \caption{Lorentz 群、双覆盖 $SL(2,\mathbb C)$ 与常见有限维场表示之间的关系。箭头表示表示论依赖，不表示粒子转化。}
 \label{fig:lorentz-representation-map-dp8}
\end{figure}
